\documentclass[11pt,a4paper]{article}
\usepackage[T1]{fontenc}
\usepackage{lmodern,amsmath,amssymb,booktabs,microtype}
\usepackage[margin=24mm]{geometry}
\usepackage[colorlinks=true,urlcolor=blue,linkcolor=blue]{hyperref}
\newcommand{\EE}{\mathbb E}
\newcommand{\tr}{\operatorname{tr}}
\newcommand{\W}[1]{\mathsf W_{#1}}
\newcommand{\Jh}{\mathcal J_h}
\newcommand{\Lh}{\mathcal L_h}
\newcommand{\Db}{\overline D}
\title{The raw memory coefficient in a central charge window\\Conserved shifts and the filtered linear field}
\author{Alan Day \and Chengheng\thanks{AI research assistant (OpenAI).}}
\date{Local research note --- 6 November 2026}
\begin{document}\maketitle
\begin{abstract}
We extend the project's second memory coefficient from exact half
filling to a compact central charge window. The five raw quartic
contractions acquire a finite charge-dependent correction entirely
along the conserved quadratic polynomial. The linear rank filter
has the same centered slope and limiting variance. Uniform rank
tails and the existing charged lifts then transfer the complete
conserved-subtracted two-energy spectrum. All finite-sector
coefficients in growing counterterms are retained.
\end{abstract}
\section{State and the central descendant table}
Use standard CAR, the overlapping complex Gaussian $p$-body
Hamiltonian, $\hbar=1$, and the normalized maximally mixed trace
in charge $m$. Norms mean $\|A\|_2^2=\EE\tr_m A^\dagger A$.
Let $h=H_m$, $\Jh A=\{h,A\}/2$, $\Lh A=[h,A]$, and
\begin{equation}
 N\text{ even},\quad m=N/2+xp,\quad |x|\leq L,\quad
 \lambda_N=p^2/N\to\lambda>0,\quad q_N=e^{-4\lambda_N}.
\end{equation}
Integer admissible charges are understood; $\lambda_N$ lies in a
positive compact set. Put $C_m=\EE\tr_mh^2\to C=J_s^2/\lambda$,
$B=h^2-C_m I$, $A=\sum_i[H,c_i]^\dagger[H,c_i]$, $O=\W2A$.
All $\W j$ are Gaussian coupling chaos, not spectral Wick products.
The output-deficit field $\mathcal S_N$ weights each quadratic
product block with input ranks $a,b$ and output rank $r$ by $S=a+b-r$.
With $D X=QX-\sum c_i^\dagger Xc_i$ and
$\mathcal T_j=(jp-D)H^j$ before restriction, define
\begin{align}
 E_0&=\{A,h^2\}/2,&M_0&=hAh,\notag\\
 S_0&=\Jh\mathcal T_3-E_0-M_0,&
 C_0&=\mathcal T_4-2E_0-M_0-2S_0,\qquad Z_0=h^4.
\end{align}
For each row $P$, the central extension of the old table is
\begin{equation}
 \W2P=C_m[p\ell_P+k_P+2\lambda_Nx\ell_P]B
       +C_m o_PO+j_P\Jh F+C_m s_P\mathcal S_N+o_{L^2}(1).
 \label{table}
\end{equation}
The old coefficients, listed on the next page, remain evaluated at
finite $\lambda_N,q_N$. In particular the extra term is conserved;
this statement concerns the central limit, not exact charge
independence at finite $N$.

\newpage
\section{Why no dynamic charge correction survives}
\begin{center}\small
\begin{tabular}{cclccc}\toprule
 $P$&$\ell_P$&$k_P$&$o_P$&$j_P$&$s_P$\\\midrule
 $E_0$&2&$-2\lambda q^2$&1&0&0\\
 $M_0$&2&$-2\lambda(2q-1)$&$q^2$&0&0\\
 $S_0$&$2q$&$-2\lambda(q^2-q+1)$&$q$&1&0\\
 $C_0$&$2q^2$&$-2\lambda(8\lambda q^2-3q^2+2q+2)$&1&$2q$&$8\lambda q^2$\\
 $Z_0$&0&$q^2+2q+3$&0&0&0\\\bottomrule
\end{tabular}
\end{center}
The exact two-charge contraction filters are those of the prior
descendant calculation~\cite{desc}. They use $n=N-m+1$,
two input ranks $a=p-u,b=p-v$, and output rank
$r=2p-u-v-S$. No charge-trace approximation is made in these filters.
For covariance eigenvalue $\kappa_j$ the central product expansion is
\begin{equation}
 \frac{\kappa_{tp-d}}{C_m}
 =e^{-4\lambda_Nt}\left[1+
 \frac{4\lambda_Nd-4\lambda_N^2t^2}{p}
 +O_{L,\lambda}\!\left(\frac{(1+d)^2}{p^2}\right)\right],
 \quad t=1,2.\label{kap}
\end{equation}
It follows by summing the logarithms of
$\prod_{j=0}^{p-1}(1-\nu/(m-j))(1-\nu/(N-m+p-j))$.
The $x$ terms cancel at order $p^{-1}$ in this ratio. Substitution
in the exact rational filters gives the five finite $B$ shifts
\begin{equation}
 4\lambda x\,(1,1,q,q^2,0)=2\lambda x(\ell_E,\ell_M,\ell_S,\ell_C,\ell_Z).
\end{equation}
All $O$, input-deficit and output-deficit coefficients agree with
the half-filled values. For the raw core
$\mathcal E=C_0-2S_0+(E_0+M_0)/2$, the net shift is
$4\lambda x(1-q)^2C_mB$, consistent with its independent filter.

Here are the norm controls promoting the expansions to
\eqref{table}. In the whole allowed rank domain the rational
denominators have size bounded below by a constant times $N$ or
$N^2$. The $A$ filters satisfy
\begin{equation}
 |g_{A,P}/C_m-o_P|\leq K(1+u+v+S)/p,\qquad |g_{A,P}/C_m|\leq K.
\end{equation}
The established central weighted $A$ norm and sixth output moment
give a squared error at most
$KC_m^4(p^{-2}+R^2/p^2+p^4/R^6)$; set $R=p^{3/4}$.
For the $H^2$ filter, subtraction of $p\ell_P$ leaves a bound
$K(1+u+v+S)$. The factorial input tail and central sixth output
moment give uniform integrability of its square. Equation
\eqref{kap} on fixed windows, followed by removal of the tails,
therefore proves strong convergence. The conversion of its
input-deficit part to $\Jh F$ is proved next.

\newpage
\section{The central linear filter and its covariance bridge}
Write $h=\sum_{a=0}^p h_a$ with
$\EE\tr_m h_a^2=C_m\omega_a$. The existing exact first-chaos
coefficient $t_a$ of $T_3=\mathcal T_3-3\Jh A$ gives
\begin{equation}
 F=\sum_a(t_a-\bar t)h_a,\qquad \bar t=\sum_a\omega_a t_a.
\end{equation}
For completeness set $k_a=\kappa_a/C_m$, $A_a=m-a$,
$B_a=N-m-a+1$, and $\alpha_a=a(N-a-p+1)/(N-2a+1)$.
The entire rank-dependent part of $t_a/C_m$ is $k_a E_a$, where
\begin{equation}
 E_a=(m-p)(1-p/A_a)(1+p/B_a)-m
 +3[\alpha_a(1+p/B_a)+(p-\alpha_a)(1-p/A_a)].
\end{equation}
Thus the exact formula, rather than a trace transfer multiplied
by $p$, gives for each fixed deficit $u$
\begin{equation}
 (t_{p-u}-t_p)/C_m\to8\lambda q u.\label{slope}
\end{equation}
Uniformly on $0\leq a\leq p$, $|E_a|\leq Kp$,
$|\partial_a E_a|\leq K$, $0\leq k_a\leq1$,
and $|\partial_a k_a|\leq K/p$. Hence the left side of
\eqref{slope} is bounded by $Ku$ over the entire rank range.
The exact probability recurrence has ratio at most $K/(u+1)$
and tends to $\lambda/(u+1)$. Consequently
$\omega_{p-u}\Rightarrow\operatorname{Pois}(\lambda)$ with
uniform factorial tails, and centering is justified in every
fixed polynomially weighted $L^2$ norm. In particular
\begin{align}
 F-8\lambda_Nq_NC_m(D_{\rm in}-\bar u h)&\to0\quad(L^2),\\
 \|F\|_2^2&\to64\lambda^3C^3q^2=64J_s^6q^2,\label{VF}\\
 \Jh F-4\lambda_Nq_NC_m L_{\rm in}&\to0\quad(L^2).
\end{align}
Here $D_{\rm in}=\sum u h_{p-u}$,
$\bar u=\sum u\omega_{p-u}$ and
$L_{\rm in}=\W2\{h,D_{\rm in}-\bar u h\}$.
The final multiplication is controlled by Gaussian fourth moments:
for an invariant linear Gaussian error $X$, $\EE X^2=vI$ and
$\EE\tr_m X^4\leq3v^2$. Thus $\|hX\|_2\leq\sqrt{3C_m v}$.

The exact body-shift ratio is
\begin{equation}
 R_j(u)=\frac{(u)_{\downarrow j}(N-2p+u+1)_{\uparrow j}
 (N-m+p-j+1)_{\uparrow j}}
 {(m-p+1)_{\uparrow j}[(p)_{\downarrow j}]^2}
 \longrightarrow\frac{(u)_{\downarrow j}}{\lambda^j}.
\end{equation}
It has a uniform polynomial bound for fixed $j$. Thus the old
signed $p,p-1,p-2$ covariance bridge and its whole-graph
total-variation estimate apply uniformly in the central window.
In particular $F/\|F\|_2$ is strongly approximated by
$\sqrt{\lambda_N}\sum_a[R_1(p-a)-1]h_a/\sqrt{C_m}$.

\newpage
\section{Charged marks and the full spectral interface}
The lift of a $(p-1)$-body field $U$ is $g(Q)U$, with
$d=m-p+1$ and $g(n)=(n-p+1)/d$. The reduced-rank binomial
identity and the direct CAR identities are valid at general charge:
\begin{align}
 \mathcal M_+(gU,V)&=(1+d^{-1})\mathcal M_+(U,V)-d^{-1}U\Db V,\\
 \mathcal M_-(gU,V)&=(1-d^{-1})\mathcal M_-(U,V)+d^{-1}UDV,\label{lift}\\
 [(kp-D)(gB')]_m&=(1-d^{-1})[(kp-1-D)B']_m
                       +(k-1)pB'_m/d.
\end{align}
Here $\mathcal M_-=-\sum[U,c_i^\dagger][V,c_i]$,
$\mathcal M_+=-\sum[U,c_i][V,c_i^\dagger]$, and $B'$ contains
one shifted slot in a fixed-length neutral product.
Reversing the order in \eqref{lift} puts $DV$ or $\Db V$ on
the left. These formulas were already derived and checked in
the mixed-spectrum note~\cite{mixed}; no new lift is postulated.

The preceding central Casimir identity gives, for a $p$-body slot,
\begin{equation}
 DV=pV,\qquad \Db V=pV-2D_{\rm in}V-E_{\rm in}V,\qquad
 (E_{\rm in}V)_a=\frac{u(u-1+N-2m)}{m-p+1}V_a.
\end{equation}
Factorial tails bound the input terms in every fixed moment.
Since $N-2m=-2xp$, the last term is $O(p^{-1})$,
and all lift corrections are $O(p/N)=O(p^{-1})$ in the
fixed-degree mixed traces. The quartic identity also has this size.

The existing oriented chord rule permits unequal index-set
sizes~\cite{BNR}; changing $p$ to $p-1$ preserves its limiting
overlap parameters. The central marked-graph trace theorem
and exponential overlap bounds transfer its fixed derivatives.
Every graph in a cross with one $F$ therefore has two signed
terms with the same limit, and cancels. The lower/upper force
and quartic lifts cover $F$--$Y_-$, $F$--$Y_+$, $F$--$R$ and
$F$--$Q_4$. The existing endpoint certificates transfer the
remaining crosses. In applying them, one first combines the raw
words into bounded marked Wick blocks; one never multiplies
an $O(p^{-1})$ trace error by an unremoved $p$ counterterm.

We obtain the same diagonal two-energy basis $(Y,F,R,T,Q_4)$
as at half filling, with parameters $q,C$. In particular
\begin{equation}
 d\nu_F=64\lambda^3C^3q^2\,d\mu_q,\qquad
 \nu_{F,Y_\pm}=\nu_{F,R}=\nu_{F,Q_4}=0.
\end{equation}
For marked representatives these follow from all fixed energy
moments, tightness, and compact moment uniqueness. Strong
replacements transfer bounded spectral tests by Cauchy--Schwarz.
Energy-weighted uses additionally employ the factorial weighted
errors and fixed-degree Gaussian/marked moment bounds above.
Ordinary chaos orthogonality alone would not give these statements.

\newpage
\section{The conserved subtraction and restored raw coefficient}
Put $d_3=q^2-q+1$, $d_5=1+q+q^2+q^3+q^4$, and
\begin{align}
 \alpha&=-2q^3/(d_3d_5),&\beta&=-q^2(1+q+q^2)/(d_3d_5),&
 \gamma&=2q/d_3,\notag\\
 \delta&=\frac{2\lambda(1-q)^3(1+q)(1+q+q^2)}{(1+q^2)d_5},&
 g_J&=\alpha+\beta+\gamma-1.&
\end{align}
Let $\mathcal Q=C_0-\alpha E_0-\beta M_0-\gamma S_0-\delta Z_0$,
$Q_4=\W4\mathcal Q$, $e_{\mathcal Q}=\EE\tr_m\mathcal Q$.
For any table column set $f_Q=f_C-\alpha f_E-\beta f_M-
\gamma f_S-\delta f_Z$. Its central correction is exactly
$2\lambda_Nx\ell_Q C_mB$ to the order under consideration.

Use actual finite-sector $a_m=\EE\tr_m A$,
$c_m=\langle O,B\rangle/\langle B,B\rangle$,
$K_A=a_m I+c_mB$, $Y=A-K_A$. Define
\begin{align}
 a&=-(1-q)^2/(1+q^2),&
 b&=\frac{2\lambda(1-q)(1+q^3)}{(1+q^2)(1+q+q^2)},\notag\\
 P_m&=(ac_m+b)h^3+[\bar t-(ac_m+b)\bar z]h,
 &\bar z&=\langle h,\W1h^3\rangle/C_m.
\end{align}
The exact cubic reduction has the same central strong remainder:
$T_3-P_m=R+a\Jh Y+F+o_{L^2}(1)$.
Here $R=\W3(T_3-a\Jh A-bh^3)$.
Indeed $B_Z,B_S=O_{L^2}(p^{-1})$ follow from the centered
$k_a$ and $k_aJ_a$ slopes, as in the preceding cubic calculation.
Their limits are $pB_Z-F/2\to0$ and
$pB_S-(1-2\lambda)F/2\to0$, also with fixed energy weights.

With $r_N=p/m$, $c=2\lambda(1-q)/(1+q)$, set
\begin{align}
 D_{2,m}&=2r_NhP_m+3r_N^2h^2K_A-K_A^2/m-mr_N^4h^4,\notag\\
 \mathcal C_{2,m}&=D_{2,m}+e_{\mathcal Q}I+\delta h^4
                 +g_Jh^2K_A+b_RhP_m\notag\\
 &\quad+C_m[p\ell_Q+k_Q+2\lambda_Nx\ell_Q+o_Qc_m
                           +8\lambda_Nq_N^2c]B.\label{counter}
\end{align}
Every term commutes with $h$ sample by sample. All $q,\lambda$
in this finite-$N$ definition mean $q_N,\lambda_N$; $C_m,c_m,
\bar t,\bar z,e_{\mathcal Q}$ remain exact sector quantities.
The renewal identity and \eqref{table} give
\begin{equation}
 U_{2,m}:=mZ_2-\mathcal C_{2,m}
 =X+b_R\Jh R+b_F\Jh F+b_J\Jh^2Y+b_D\Lh^2Y+C\kappa Y+o_{L^2}(1),
\end{equation}
Here $\widehat M(s)=G+isZ_1-s^2Z_2/2+O(s^3)$ defines the
finite-$N$ centered memory coefficients. Also
$T=(1-\Pi_N)\mathcal S_N$ with
$\Pi_N=\operatorname{proj}_{\{B,L_{\rm in},O\}}$,
$X=Q_4+8\lambda Cq^2T$, and
\begin{align}
 b_R&=-2(1-q)^2/d_3,& b_F&=2(q-1),& b_D&=q^2/(4d_5),\notag\\
 b_J&=g_J+ab_R,&\kappa&=(1-q)(1+q)^2/d_5+8\lambda q^2\rho,
\end{align}
$\rho=v/(v+\Delta)$, $v=2\lambda(1-q)+16\lambda^2q/(1+q)$,
$\Delta=2\lambda(1+q)$. The renewal error is $O_{L^2}(p^{-1})$:
its terms have factors $r_N$, $r_N^2$, or $a_m/m=O(p^{-1})$
and bounded centered higher moments. The finite conserved
$K_A^2/m$ term is retained in \eqref{counter}.

\newpage
\section{Spectrum, evidence, and continuation}
Let $e=(E+E')/2$, $\omega=E-E'$,
$H_2=b_Je^2+b_D\omega^2+C\kappa$. The result is
\begin{equation}
 d\nu_{U_2}=V_4d\mu_{q^4}+V_2d\mu_{q^2}
 +b_R^2e^2V_Rd\mu_{q^3}+b_F^2e^2V_Fd\mu_q
 +WH_2^2d\mu_{q^2}.\label{spectrum}
\end{equation}
These are the same explicit weights as in the raw-spectrum
note~\cite{raw}, with $W=4C^2\lambda+16C^2\lambda^2q/(1+q)$,
$V_F$ in \eqref{VF}, and $V_2=64\lambda^2C^4q^4v\Delta/(v+\Delta)$.
The existing Mehler measure $\mu_r$ has variance-$C$ marginals.
The cross measures are
\begin{equation}
 d\nu_{Y,U_2}=WH_2d\mu_{q^2},\qquad
 d\nu_{L,U_2}=aeWH_2d\mu_{q^2}+b_FeV_Fd\mu_q+b_ReV_Rd\mu_{q^3}.
\end{equation}
For densities $w=WK_2$, $f=V_FK_1$, $r=V_RK_3$,
$z=V_4K_4+V_2K_2$, the $(Y,L,U_2)$ matrix is
\begin{equation}
 M\operatorname{diag}(w,f,r,z)M^T,\qquad
 M=\begin{pmatrix}1&0&0&0\\ae&1&1&0\\H_2&b_Fe&b_Re&1\end{pmatrix}.
\end{equation}
Its determinant is $w[(f+r)z+e^2fr(b_F-b_R)^2]$.
Thus the compact central window changes the required subtraction,
but not the limiting dynamic spectral matrix. The ordered
small-$\lambda$ shape already derived in Ref.~\cite{raw} is
inherited after this central double-scaled limit.

\paragraph{Executed evidence.}
The restored symbolic certificate checks all five central filters,
the cubic slope, and the one-body-shift ratio. Its recorded source
hashes were verified before reuse. A new calculation checks
123 exact rational coefficient references and 36 full-rank sums
at $\lambda=0.1,0.25,0.5$, $p=16,32,64,128$, $x=-2,0,2$,
each at 80 and 120 decimal digits. All $p+1$ ranks are included;
finite variance differences, signed-covariance errors, weighted
tails, and both cubic remainders are saved. Agreement of numerical
precisions is not a convergence proof or an interval certificate.
The assembly verifies seven prior/current certificates and the
positive matrix algebra, without repeating old endpoint scans.

\paragraph{Next calculation.}
The next issue is charge summation: obtain bounds outside every
fixed central window for the \emph{conserved-subtracted} fields,
then combine them with the binomial charge law. Compact-$x$
uniformity alone cannot interchange that sum with the limit.
Afterwards control the two-insertion finite-memory remainder and
connect the bilocal to the Susskind bulk meson dictionary. This
note identifies fixed Taylor coefficients, not the finite-memory kernel.

\begin{thebibliography}{9}\small
\bibitem{BNR} M. Berkooz, V. Narovlansky and H. Raj,
\emph{Complex Sachdev--Ye--Kitaev model in the double scaling limit},
\href{https://arxiv.org/html/2006.13983v2}{arXiv:2006.13983v2}, Sec.~3.3.
\bibitem{desc} A. Day and Chengheng, local notes:
\emph{Restoring the raw memory descendants};
\emph{Central coherent norms}; \emph{Central force spectrum} (2026).
\bibitem{mixed} A. Day and Chengheng, local notes:
\emph{Mixed spectra of the first memory fields};
\emph{Linear rank-filter covariance bridge} (2026).
\bibitem{raw} A. Day and Chengheng,
\emph{The second coefficient after a conserved subtraction},
local six-page note, 6 November 2026.
\end{thebibliography}
\end{document}
